% !TeX spellcheck = en_US
\addsection[2]{Rule changes for competitive play}{\art/haste.png}

\iftoggle{printable}{}{\bigbreak}

\begin{multicols*}{2}
We've compiled some custom rules and best practices to enhance your gaming experience.
Keep in mind that these are just the \textbf{opinions of the authors} of this book.
Enjoy!

\subsection*{Magic Arrows}

Except for the Magic Arrow, each Spell appears exactly twice in the Spell Deck.
After players prepare their starting M\&M Decks, it is recommended to remove Magic Arrows from the Spell Deck until no more than two remain.

\subsection*{Guaranteed Settlement}

When distributing Far Tiles (II--III) between players, each player's pool of Map Tiles must include one Settlement.

\subsection*{Round One Mulligan}

At the start of the \nth{1} Round, each player may reshuffle their hand of Cards back into their Deck of Might and Magic and draw a new starting hand of Cards once.

\subsection*{Tier V--VII Combat on Hard$+$\\Difficulty}

Treat Level V--VII Combats as one difficulty step lower when playing on Hard or Impossible.
For example, a Level VII Combat on Hard will be against 2 × \svg{azure} Neutral Units.

{
  \hfill
  \transparent{0.2}{\includegraphics[width=0.7\linewidth]{\art/black_market.png}}
  \vspace{-3em}
}

\subsection*{Trading Post}

Although rules as written state that you can remove only one Card from your hand while Visiting a Trading Post, it is recommended to allow removing as many as you like.

\subsection*{Victory Points}

When playing with the optional Victory Points from the Tournament book, the following rules apply (unless overruled by the Scenario):
\begin{itemize}
  \item Faction Towns grant 1 VP.
  \item Possessing the Grail at the end of the match grants 3 VPs. If the Grail was transported to that player's Faction Town, it grants 5 VPs instead.
  \item Capturing a Dragon Utopia grants 3 VPs.
\end{itemize}

\end{multicols*}

\begin{multicols}{2}

\subsection*{Combat with Neutral Units}

During Combat with Neutral Units in Clash Scenarios, if a player decides to deploy their Units only on one half (left or right) of the Combat Board, the player positioning the Neutral Units must do the same.

Placing Neutral Units far apart only makes the other player waste \svg{movement}, which slows the game down.
The exception to this rule is when there are three Neutral \svg{unit_ground} or \svg{unit_flying} Units to position, as all of them must go to the front line.
See examples that follow.

\begin{center}
  \transparent{0.2}{\includegraphics[width=0.65\linewidth]{\art/resurrection.png}}
\end{center}

\end{multicols}

\begin{figure*}[!h]
  \mbox{}
  \hfill
  \begin{minipage}[t]{0.44\textwidth}
    \centering
    \includegraphics[width=\linewidth]{\examples/battle_good.png}
    \caption[good protected]{\textit{Neutral Units are positioned correctly.}}
  \end{minipage}
  \hfill
  \begin{minipage}[t]{0.44\textwidth}
    \centering
    \includegraphics[width=\linewidth]{\examples/battle_bad.png}
    \caption[bad protected]{\textit{This deployment is not allowed because the Necropolis player placed their Units on the left half of the Combat Board.
      The Peasants must also start the Combat on the left side.}}
  \end{minipage}
  \hfill
  \mbox{}
\end{figure*}

\clearpage

\begin{figure*}[!h]
  \mbox{}
  \hfill
  \begin{minipage}[t]{0.44\textwidth}
    \centering
    \includegraphics[width=\linewidth]{\examples/battle_exception.png}
    \caption[exception protected]{\textit{\textbf{Exception:} Three non-\svg{unit_ranged} Units must be positioned on the front line.}}
  \end{minipage}
  \hfill
  \begin{minipage}[t]{0.44\textwidth}
    \centering
    \includegraphics[width=\linewidth]{\examples/battle_ranged_bad.png}
    \caption[ranged protected]{\textit{The Boars must occupy one of the green Fields.}}
  \end{minipage}
  \hfill
  \mbox{}
\end{figure*}

\bigskip
\bigskip
\subheader{Rebalanced cards}

\begin{table*}[!hb]
  \hommtable{50}{
    \centering
    \begin{tabularx}{0.95\linewidth}{
        >{\centering}p{0.15\linewidth}
        >{\centering\arraybackslash}X
        >{\centering\arraybackslash}X
      }
      \darkcell[1.2]{Ability}
      & \darkcell[1.2]{Current text}
      & \darkcell[1.2]{Revised text} \\

      % Artillery
      \darkcell[3.8]{Artillery\\[0.3em] \includegraphics[width=0.8\linewidth]{\abilities/artillery.png}}
      & \lightcell[3.8]{\svg{instant-yellow} Deal 1 damage to an enemy unit with the lowest initiative.\\[0.3em]
        --- \svg{expert-yellow} ---\\[0.3em]
        \svg{instant-yellow} When using the Ballista card, resolve its effect against the same target 3 times.}
      & \lightcell[3.8]{\svg{activation-yellow} Deal 1 damage to an enemy unit.\\[0.3em]
        --- \svg{expert-yellow} ---\\[0.3em]
        \svg{instant-yellow} When using the Ballista card, resolve its effect against the same target 3 times. You may select the target.} \\

      % Ballistics
      \darkcell[5.0]{Ballistics\\[0.3em] \includegraphics[width=0.8\linewidth]{\abilities/ballistics.png}}
      & \lightcell[5.0]{\svg{instant-yellow} Play this card during a siege.\\Destroy 1 Wall or Gate.\\[0.3em]
        --- \svg{expert-yellow} ---\\[0.3em]
        \svg{instant-yellow} Play this card during a siege.\\Destroy the Arrow Tower.}
      & \lightcell[5.0]{\svg{instant-yellow} Play this card during a siege.\\Destroy 1 Wall, Gate or Arrow Tower.\\[0.3em]
        --- \svg{expert-yellow} ---\\[0.3em]
        \svg{instant-yellow} Play this card during a siege.\\Destroy all Walls, Gate, and Arrow Tower.\\Every unit adjacent to a Wall or Gate obstacle suffers 1 damage.} \\

      % Estates
      \darkcell[3.0]{Estates\\[0.3em] 
\includegraphics[width=0.8\linewidth]{\abilities/estates.png}}
      & \lightcell[3.0]{\svg{instant-yellow} Gain 3 Gold.\\[0.3em]
        --- \svg{expert-yellow} ---\\[0.3em]
        \svg{instant-yellow} Gain 6 Gold.}
      & \lightcellnerf[3.0]{\svg{instant-yellow} Gain 2 Gold.\\[0.3em]
        --- \svg{expert-yellow} ---\\[0.3em]
        \svg{instant-yellow} Gain 4 Gold.} \\

      % First Aid
      \darkcell[5.5]{First Aid\\[0.3em] 
\includegraphics[width=0.8\linewidth]{\abilities/first_aid.png}}
      & \lightcell[5.5]{\svg{instant-yellow} Remove 1 damage from one of your units.\\[0.3em]
        --- \svg{expert-yellow} ---\\[0.3em]
        \svg{instant-yellow} When using the First Aid Tent card, resolve its effect against the same target 3 times.}
      & \lightcell[5.5]{\svg{instant-yellow} You may remove any effect (or Paralysis) from the selected unit.\\Remove 1 damage from the selected unit.\\[0.3em]
        --- \svg{expert-yellow} ---\\[0.3em]
        \svg{instant-yellow} When using the First Aid Tent card, resolve its effect against the same target 3 times.\\You may remove any effect or (Paralysis) from the selected unit.} \\

      % Intelligence
      \darkcell[5.0]{Intelligence\\[0.3em] 
\includegraphics[width=0.8\linewidth]{\abilities/intelligence.png}}
      & \lightcell[5.0]{\svg{instant-yellow} During Combat, before any unit activates, play a Spell card.\\You can still only play a Spell card during a Combat round.\\[0.3em]
        --- \svg{expert-yellow} ---\\[0.3em]
        \svg{instant-yellow} During Combat, before any unit activates, play a Spell card.\\This spell does not count toward your spell limit per Combat round.}
      & \lightcell[5.0]{\svg{instant-yellow} Play a spell from your discard pile.\\[0.3em]
        --- \svg{expert-yellow} ---\\[0.3em]
        \svg{instant-yellow} Play a spell from your discard pile.\\This spell does not count toward your spell limit per Combat round.} \\


    \end{tabularx}
  }
\end{table*}

\begin{table*}[!hb]
  \hommtable{50}{
    \centering
    \begin{tabularx}{0.95\linewidth}{
        >{\centering}p{0.15\linewidth}
        >{\centering\arraybackslash}X
        >{\centering\arraybackslash}X
      }
      \darkcell[1.2]{Ability}
      & \darkcell[1.2]{Current text}
      & \darkcell[1.2]{Revised text} \\

      % Leadership (Expert only)
      \darkcell[3.4]{Leadership (Expert only)\\[0.3em] 
\includegraphics[width=0.8\linewidth]{\abilities/leadership.png}}
      & \lightcell[3.4]{--- \svg{expert-yellow} ---\\[0.3em]
      \svg{instant-yellow} Draw 2 cards. Then, gain a \svg{morale_positive-yellow} token.}
      & \lightcellnerf[3.4]{--- \svg{expert-yellow} ---\\[0.3em]
      \svg{instant-yellow} Draw 2 cards.} \\

      % Luck
      \darkcell[3.5]{Luck\\[0.3em] 
\includegraphics[width=0.8\linewidth]{\abilities/luck.png}}
      & \lightcell[3.5]{\svg{ongoing-yellow} You can reroll a Treasure die and a Resource die once during this turn.\\[0.3em]
        --- \svg{expert-yellow} ---\\[0.3em]
        \svg{ongoing-yellow} You can reroll any die once during this turn.}
      & \lightcell[3.5]{\svg{ongoing-yellow} This turn, you may reroll each treasure or resource die once.\\[0.3em]
        --- \svg{expert-yellow} ---\\[0.3em]
        \svg{ongoing-yellow} This turn, you may reroll each die once.} \\

      % Mysticism (Basic only)
      \darkcell[3.4]{Mysticism (Basic only)\\[0.3em] 
\includegraphics[width=0.8\linewidth]{\abilities/mysticism.png}}
      & \lightcell[3.4]{\svg{instant-yellow} Play this card immediately after casting a spell.\\Instead of discarding the Spell card, take it back into your hand.}
      & \lightcell[3.4]{\svg{instant-yellow} Play this card immediately after casting a spell.\\Instead of discarding the Spell card, take it back into your hand as well as 1 card played alongside it.} \\

      % Tactics
      \darkcell[4.5]{Tactics\\[0.3em] 
\includegraphics[width=0.8\linewidth]{\abilities/tactics.png}}
      & \lightcell[4.5]{\svg{instant-yellow} At the start of Combat, you can switch the position of any 2 of your units.\\[0.3em]
        --- \svg{expert-yellow} ---\\[0.3em]
        \svg{activation-yellow} During Combat, you can switch the position of any 2 of your units.}
      & \lightcell[4.5]{\svg{instant-yellow} At the start of combat, rearrange your units starting positions.\\[0.3em]
        --- \svg{expert-yellow} ---\\[0.3em]
        \svg{instant-yellow} At the start of combat, shuffle your current hand and discard pile back into your draw pile. Draw back up to your current hand limit.} \\

      % Wisdom
      \darkcell[6.0]{Wisdom\\[0.3em] 
\includegraphics[width=0.8\linewidth]{\abilities/wisdom.png}}
      & \lightcell[6.0]{\svg{instant-yellow} The cost of buying spells in this Town is reduced by 2 gold.\\When buying spells from your Magic Guild, do Search~(3) instead of Search~(2).\\[0.3em]
        --- \svg{expert-yellow} ---\\[0.3em]
        \svg{instant-yellow} The cost of buying spells in this Town is reduced by 2 gold.\\When buying spells from your Magic Guild, do Search~(4) instead of Search~(2).}
      & \lightcell[6.0]{\svg{instant-yellow} +1~\svg{empower-yellow}. The next Spell you cast does not count toward the limit.\\[0.3em]
        --- \svg{expert-yellow} ---\\[0.3em]
        \svg{instant-yellow} +2~\svg{empower-yellow}. The next Spell you cast does not count toward the limit.} \\

    \end{tabularx}
  }
\end{table*}
